\shead{Tìm tỉ số phần trăm của hai số}
\begin{knowbox}\begin{itemize}[leftmargin=1.4em,itemsep=3pt]
\item Muốn tìm tỉ số phần trăm của hai số $A$ và $B$, ta tìm thương của $A$ và $B$ rồi viết thương đó dưới dạng tỉ số phần trăm (nhân với $100$ và viết thêm kí hiệu $\%$).\par
Ví dụ: $\ $tỉ số phần trăm của $9$ và $12$ là $\ 9:12=0{,}75=75\%$.
\item Đọc kĩ đề bài: ``$A$ so với $B$'' thì $B$ là cả đại lượng ($100\%$) và ta chia cho $B$.
\item \textbf{Thực hiện kế hoạch:} coi kế hoạch là $100\%$; $\ \text{số thực hiện}:\text{số kế hoạch}$ cho biết đã đạt bao nhiêu phần trăm kế hoạch; phần vượt quá $100\%$ là số phần trăm \emph{vượt mức} kế hoạch.
\item \textbf{Tăng (giảm) bao nhiêu phần trăm:} $\ (\text{phần tăng hoặc giảm}):(\text{số lúc đầu})$, viết dưới dạng tỉ số phần trăm.
\end{itemize}
\end{knowbox}
\gap
\shead{Tìm $a\%$ của một số}
\begin{knowbox}\begin{itemize}[leftmargin=1.4em,itemsep=3pt]
\item Muốn tìm $a\%$ của $M$, ta lấy $\ M\times a:100\ $ hoặc $\ M:100\times a$.\par
Ví dụ: $\ 15\%$ của $60$ kg là $\ 60\times15:100=9$ (kg).
\item $M:100$ là $1\%$ của $M$; nhân với $a$ ta được $a\%$ của $M$. \ Cách ``rút về $1\%$'' này cũng giúp ta tính ngược: nếu $a\%$ của một số là $N$ thì $1\%$ của số đó là $N:a$, và số đó là $N:a\times100$.
\end{itemize}
\end{knowbox}
\gap
\shead{Tỉ số phần trăm trong đời sống}
\begin{knowbox}\begin{itemize}[leftmargin=1.4em,itemsep=3pt]
\item \textbf{Giảm giá:} giảm giá $a\%$ nghĩa là ta trả $(100-a)\%$ giá cũ.
\item \textbf{Tiền lãi:} tiền lãi $=$ giá bán $-$ giá vốn (giá mua). \ ``Lãi $20\%$'' thường hiểu là lãi $20\%$ giá vốn, khi đó giá bán bằng $120\%$ giá vốn.
\item \textbf{Lãi suất tiết kiệm:} với lãi suất $0{,}5\%$ một tháng, mỗi tháng ngân hàng trả thêm $0{,}5\%$ số tiền đang gửi.
\item \textbf{Nước muối:} ``nước muối $5\%$'' nghĩa là lượng muối bằng $5\%$ lượng nước muối.
\end{itemize}
\note{Luôn tự hỏi: ``phần trăm \emph{của cái gì}?'' -- đại lượng đó ứng với $100\%$.}
\end{knowbox}
\gap
\begin{exbox}{\ExWord\ 1}
Tìm tỉ số phần trăm của: \ \pa{a} $0{,}6$ và $2{,}4$; \ \pa{b} $7$ và $4$.
\guide
\pa{a} $0{,}6:2{,}4=0{,}25=25\%$.\par
\pa{b} $7:4=1{,}75=175\%$.
\dapso{\pa{a} $25\%$;\ \ \pa{b} $175\%$.}
\end{exbox}
\gap
\begin{exbox}{\ExWord\ 2}
Lớp 5B có $18$ học sinh nam và $22$ học sinh nữ. Tính tỉ số phần trăm của số học sinh nam, số học sinh nữ so với số học sinh cả lớp.
\guide
Số học sinh cả lớp: $18+22=40$ (học sinh).\par
Tỉ số phần trăm của số học sinh nam so với số học sinh cả lớp: $18:40=0{,}45=45\%$.\par
Tỉ số phần trăm của số học sinh nữ so với số học sinh cả lớp:\par
\emph{Cách 1.} $22:40=0{,}55=55\%$. \qquad \emph{Cách 2.} $100\%-45\%=55\%$.
\dapso{nam $45\%$; nữ $55\%$.}
\end{exbox}
\gap
\begin{exbox}{\ExWord\ 3}
Tính $35\%$ của $240$.
\guide
\emph{Cách 1.} $240:100\times35=2{,}4\times35=84$.\par
\emph{Cách 2.} $240\times35:100=8\,400:100=84$.
\dapso{$84$.}
\end{exbox}
\gap
\begin{exbox}{\ExWord\ 4}
Một chi đoàn thanh niên dự định trồng $250$ cây, nhưng thực tế đã trồng được $290$ cây. Hỏi chi đoàn đã đạt bao nhiêu phần trăm kế hoạch và vượt mức kế hoạch bao nhiêu phần trăm?
\guide
Tỉ số phần trăm của số cây đã trồng so với kế hoạch: $290:250=1{,}16=116\%$.\par
Chi đoàn đã vượt mức kế hoạch: $116\%-100\%=16\%$.
\dapso{đạt $116\%$ kế hoạch; vượt mức $16\%$.}
\end{exbox}
\gap
\begin{exbox}{\ExWord\ 5}
Một chiếc cặp sách giá $180\,000$ đồng. Trong đợt khuyến mãi, giá chiếc cặp được giảm $15\%$. Hỏi khi khuyến mãi, chiếc cặp có giá bao nhiêu tiền?
\guide
\emph{Cách 1.} Số tiền được giảm: $180\,000\times15:100=27\,000$ (đồng).\par
Giá khi khuyến mãi: $180\,000-27\,000=153\,000$ (đồng).\par
\emph{Cách 2.} Người mua trả $100\%-15\%=85\%$ giá cũ: $\ 180\,000\times85:100=153\,000$ (đồng).
\dapso{$153\,000$ đồng.}
\end{exbox}
\dividerbar{Luyện tập}
\begin{probbox}{1}{\relax}
Tìm tỉ số phần trăm của hai số.
\begin{cols3}
\pa{a} $18$ và $24$ & \pa{b} $4{,}5$ và $12{,}5$ & \pa{c} $0{,}9$ và $3{,}6$\\
\pa{d} $27$ và $15$ & \pa{e} $3{,}4$ và $8$ & \pa{g} $\dfrac14$ và $5$
\end{cols3}
\end{probbox}
\probsol{
\pa{a} $18:24=0{,}75=75\%$. \quad \pa{b} $4{,}5:12{,}5=0{,}36=36\%$. \quad \pa{c} $0{,}9:3{,}6=0{,}25=25\%$.\par
\pa{d} $27:15=1{,}8=180\%$. \quad \pa{e} $3{,}4:8=0{,}425=42{,}5\%$. \quad \pa{g} $\dfrac14:5=0{,}05=5\%$.
\dapso{$75\%$; $36\%$; $25\%$; $180\%$; $42{,}5\%$; $5\%$.}}
\gap
\begin{probbox}{2}{\relax}
Tính.
\begin{cols3}
\pa{a} $15\%$ của $60$ kg & \pa{b} $8\%$ của $450$ m & \pa{c} $25\%$ của $7{,}2$ l\\
\pa{d} $40\%$ của $3{,}5$ tấn & \pa{e} $120\%$ của $85$ cm & \pa{g} $12{,}5\%$ của $96$ m$^2$
\end{cols3}
\end{probbox}
\probsol{
\pa{a} $60\times15:100=9$ (kg). \quad \pa{b} $450\times8:100=36$ (m). \quad \pa{c} $7{,}2\times25:100=1{,}8$ (l).\par
\pa{d} $3{,}5\times40:100=1{,}4$ (tấn). \quad \pa{e} $85\times120:100=102$ (cm).\par
\pa{g} $96\times12{,}5:100=12$ (m$^2$).
\dapso{$9$ kg; $36$ m; $1{,}8$ l; $1{,}4$ tấn; $102$ cm; $12$ m$^2$.}}
\gap
\begin{probbox}{3}{\relax}
Một túi có $48$ viên bi, trong đó có $30$ viên bi đỏ, còn lại là bi xanh. \ \pa{a} Số bi đỏ chiếm bao nhiêu phần trăm số bi trong túi? \ \pa{b} Tính tỉ số phần trăm của số bi xanh so với số bi đỏ.
\end{probbox}
\probsol{
\pa{a} $30:48=0{,}625=62{,}5\%$.\par
\pa{b} Số bi xanh: $48-30=18$ (viên). \ $18:30=0{,}6=60\%$.
\dapso{\pa{a} $62{,}5\%$;\ \ \pa{b} $60\%$.}}
\gap
\begin{probbox}{4}{\relax}
Trong bài kiểm tra Toán, lớp 5C có $34$ trên $40$ học sinh đạt điểm tốt. Hỏi số học sinh đạt điểm tốt chiếm bao nhiêu phần trăm số học sinh cả lớp?
\end{probbox}
\probsol{
$34:40=0{,}85=85\%$.
\dapso{$85\%$.}}
\gap
\begin{probbox}{5}{\relax}
Một xưởng xe đạp dự định làm $400$ chiếc xe trong một tháng, nhưng thực tế đã làm được $460$ chiếc. Hỏi xưởng đã đạt bao nhiêu phần trăm kế hoạch và vượt mức kế hoạch bao nhiêu phần trăm?
\end{probbox}
\probsol{
$460:400=1{,}15=115\%$.\par
Vượt mức kế hoạch: $115\%-100\%=15\%$.
\dapso{đạt $115\%$; vượt mức $15\%$.}}
\gap
\begin{probbox}{6}{\relax}
Năm nay một xã dự định trồng $1\,250$ cây. Đến hết tháng Ba, xã đã trồng được $1\,000$ cây. Hỏi xã đã thực hiện được bao nhiêu phần trăm kế hoạch và còn phải thực hiện bao nhiêu phần trăm nữa?
\end{probbox}
\probsol{
Đã thực hiện: $1\,000:1\,250=0{,}8=80\%$ (kế hoạch).\par
Còn phải thực hiện: $100\%-80\%=20\%$ (kế hoạch).
\dapso{đã thực hiện $80\%$; còn $20\%$.}}
\gap
\begin{probbox}{7}{\relax}
Mỗi tháng cô Lan được lĩnh $12$ triệu đồng tiền lương và để dành được $2{,}7$ triệu đồng. Hỏi số tiền để dành chiếm bao nhiêu phần trăm tiền lương?
\end{probbox}
\probsol{
$2{,}7:12=0{,}225=22{,}5\%$.
\dapso{$22{,}5\%$.}}
\gap
\begin{probbox}{8}{\relax}
Năm ngoái Minh cân nặng $38$ kg. Năm nay cân nặng của Minh tăng $10\%$ so với năm ngoái. Tính cân nặng của Minh năm nay.
\end{probbox}
\probsol{
Số ki-lô-gam tăng thêm: $38\times10:100=3{,}8$ (kg).\par
Cân nặng năm nay: $38+3{,}8=41{,}8$ (kg).
\dapso{$41{,}8$ kg.}}
\gap
\begin{probbox}{9}{\relax}
Một mảnh đất hình chữ nhật có chiều dài $50$ m, chiều rộng bằng $60\%$ chiều dài. Người ta dùng $45\%$ diện tích để trồng rau, $35\%$ diện tích để trồng hoa, phần còn lại làm lối đi. Tính diện tích lối đi.
\end{probbox}
\probsol{
Chiều rộng: $50\times60:100=30$ (m). \ Diện tích mảnh đất: $50\times30=1\,500$ (m$^2$).\par
Lối đi chiếm: $100\%-45\%-35\%=20\%$ (diện tích).\par
Diện tích lối đi: $1\,500\times20:100=300$ (m$^2$).
\dapso{$300$ m$^2$.}}
\gap
\begin{probbox}{10}{\relax}
Một cửa hàng bán được $1\,200$ kg gạo, trong đó $35\%$ là gạo nếp. Hỏi cửa hàng bán được bao nhiêu ki-lô-gam gạo tẻ?
\end{probbox}
\probsol{
Gạo tẻ chiếm: $100\%-35\%=65\%$ (số gạo).\par
$1\,200\times65:100=780$ (kg).
\dapso{$780$ kg.}}
\gap
\begin{probbox}{11}{\relax}
Một trường có $640$ học sinh, trong đó $55\%$ là học sinh nữ. Hỏi trường có bao nhiêu học sinh nam?
\end{probbox}
\probsol{
Số học sinh nữ: $640\times55:100=352$ (học sinh).\par
Số học sinh nam: $640-352=288$ (học sinh).
\dapso{$288$ học sinh nam.}}
\gap
\begin{probbox}{12}{\relax}
Một cô bán hàng mua một chiếc áo với giá $150\,000$ đồng và bán ra với giá $180\,000$ đồng. Hỏi tiền lãi bằng bao nhiêu phần trăm giá mua?
\end{probbox}
\probsol{
Tiền lãi: $180\,000-150\,000=30\,000$ (đồng).\par
$30\,000:150\,000=0{,}2=20\%$.
\dapso{$20\%$ giá mua.}}
\gap
\begin{probbox}{13}{\relax}
Một quyển truyện giá $60\,000$ đồng. Tại hội sách, quyển truyện được bán giảm giá $10\%$. Hỏi tại hội sách, quyển truyện có giá bao nhiêu tiền?
\end{probbox}
\probsol{
Người mua trả: $100\%-10\%=90\%$ (giá bìa).\par
$60\,000\times90:100=54\,000$ (đồng).
\dapso{$54\,000$ đồng.}}
\gap
\begin{probbox}{14}{\relax}
Bà Hạnh gửi tiết kiệm $20\,000\,000$ đồng với lãi suất $0{,}6\%$ một tháng. Hỏi sau một tháng, bà Hạnh nhận được bao nhiêu tiền lãi và có tất cả bao nhiêu tiền trong sổ tiết kiệm?
\end{probbox}
\probsol{
Tiền lãi: $20\,000\,000\times0{,}6:100=120\,000$ (đồng).\par
Số tiền trong sổ: $20\,000\,000+120\,000=20\,120\,000$ (đồng).
\dapso{tiền lãi $120\,000$ đồng; tất cả $20\,120\,000$ đồng.}}
\gap
\begin{probbox}{15}{\relax}
Nước biển chứa $3{,}5\%$ muối. Hỏi trong $40$ kg nước biển có bao nhiêu ki-lô-gam muối?
\end{probbox}
\probsol{
$40\times3{,}5:100=1{,}4$ (kg).
\dapso{$1{,}4$ kg.}}
\gap
\begin{probbox}{16}{\relax}
Giá một quyển vở tăng từ $8\,000$ đồng lên $9\,000$ đồng. Hỏi giá quyển vở đã tăng bao nhiêu phần trăm?
\end{probbox}
\probsol{
Số tiền tăng thêm: $9\,000-8\,000=1\,000$ (đồng).\par
So với giá cũ: $1\,000:8\,000=0{,}125=12{,}5\%$.
\dapso{$12{,}5\%$.}}
\gap
\begin{probbox}{17}{\relax}
Giá một đôi giày thể thao giảm từ $250\,000$ đồng xuống còn $200\,000$ đồng. Hỏi giá đôi giày đã giảm bao nhiêu phần trăm?
\end{probbox}
\probsol{
Số tiền giảm: $250\,000-200\,000=50\,000$ (đồng).\par
So với giá cũ: $50\,000:250\,000=0{,}2=20\%$.
\dapso{$20\%$.}}
\gap
\begin{probbox}{18}{\relax}
Một đội bóng đá đã thi đấu $25$ trận: thắng $16$ trận, hoà $4$ trận, còn lại là thua. Tính tỉ số phần trăm số trận thắng, hoà, thua so với tổng số trận.
\end{probbox}
\probsol{
Số trận thua: $25-16-4=5$ (trận).\par
Thắng: $16:25=0{,}64=64\%$; \ hoà: $4:25=0{,}16=16\%$; \ thua: $5:25=0{,}2=20\%$.
\dapso{thắng $64\%$; hoà $16\%$; thua $20\%$.}}
\gap
\begin{probbox}{19}{\relax}
Hoa đọc một quyển sách dày $240$ trang. Ngày thứ nhất Hoa đọc được $30\%$ số trang, ngày thứ hai đọc được $45\%$ số trang. Hỏi còn bao nhiêu trang chưa đọc?
\end{probbox}
\probsol{
Số trang chưa đọc chiếm: $100\%-30\%-45\%=25\%$ (số trang).\par
$240\times25:100=60$ (trang).
\dapso{$60$ trang.}}
\gap
\begin{probbox}{20}{\relax}
Một bác nông dân trồng $450$ cây giống, có $18$ cây bị chết. Hỏi số cây sống chiếm bao nhiêu phần trăm số cây đã trồng?
\end{probbox}
\probsol{
Số cây sống: $450-18=432$ (cây).\par
$432:450=0{,}96=96\%$.
\dapso{$96\%$.}}
\gap
\begin{probbox}{21}{\relax}
Trong $500$ hạt đỗ đem gieo, có $92\%$ số hạt nảy mầm. Hỏi có bao nhiêu hạt không nảy mầm?
\end{probbox}
\probsol{
Số hạt nảy mầm: $500\times92:100=460$ (hạt).\par
Số hạt không nảy mầm: $500-460=40$ (hạt).
\dapso{$40$ hạt.}}
\gap
\begin{probbox}{22}{\relax}
Lớp 5A có $35$ học sinh, trong đó $40\%$ số học sinh đạt danh hiệu Học sinh Xuất sắc. Hỏi lớp 5A có bao nhiêu học sinh không đạt danh hiệu đó?
\end{probbox}
\probsol{
Số học sinh xuất sắc: $35\times40:100=14$ (học sinh).\par
Số học sinh không đạt: $35-14=21$ (học sinh).
\dapso{$21$ học sinh.}}
\gap
\begin{probbox}{23}{\relax}
Một giỏ có $80$ quả táo và lê, trong đó số lê chiếm $40\%$. Tính tỉ số phần trăm của số táo so với số lê.
\end{probbox}
\probsol{
Số lê: $80\times40:100=32$ (quả). \ Số táo: $80-32=48$ (quả).\par
Tỉ số phần trăm của số táo so với số lê: $48:32=1{,}5=150\%$.
\dapso{$150\%$.}}
\gap
\begin{probbox}{24}{\relax}
Trong ngày hội trồng cây, lớp 5A có $32$ trên $40$ học sinh tham gia, lớp 5B có $27$ trên $30$ học sinh tham gia. Lớp nào có tỉ lệ học sinh tham gia cao hơn?
\end{probbox}
\probsol{
Lớp 5A: $32:40=0{,}8=80\%$.\par
Lớp 5B: $27:30=0{,}9=90\%$.\par
$90\%>80\%$ nên lớp 5B có tỉ lệ tham gia cao hơn.
\dapso{lớp 5B ($90\%$ so với $80\%$).}}
\gap
\begin{probbox}{25}{\relax}
Năm ngoái một nhà máy sản xuất được $2\,400$ sản phẩm. Năm nay nhà máy sản xuất nhiều hơn năm ngoái $15\%$. Hỏi năm nay nhà máy sản xuất được bao nhiêu sản phẩm?
\end{probbox}
\probsol{
Số sản phẩm năm nay bằng: $100\%+15\%=115\%$ (số sản phẩm năm ngoái).\par
$2\,400\times115:100=2\,760$ (sản phẩm).
\dapso{$2\,760$ sản phẩm.}}
\gap
\begin{probbox}{26}{\relax}
Vé vào bảo tàng giá $40\,000$ đồng một vé. Trẻ em được giảm $25\%$ giá vé. Hỏi một gia đình gồm $2$ người lớn và $3$ trẻ em phải trả bao nhiêu tiền vé?
\end{probbox}
\probsol{
Giá vé trẻ em: $40\,000\times(100-25):100=30\,000$ (đồng).\par
Gia đình phải trả: $40\,000\times2+30\,000\times3=80\,000+90\,000=170\,000$ (đồng).
\dapso{$170\,000$ đồng.}}
\gap
\begin{probbox}{27}{\relax}
Một mảnh vườn hình chữ nhật có chiều dài $25$ m, chiều rộng $16$ m. Người ta tăng chiều dài thêm $20\%$, giữ nguyên chiều rộng. Tính diện tích mới của mảnh vườn. Diện tích đã tăng thêm bao nhiêu phần trăm?
\end{probbox}
\probsol{
Chiều dài mới: $25\times120:100=30$ (m). \ Diện tích mới: $30\times16=480$ (m$^2$).\par
Diện tích cũ: $25\times16=400$ (m$^2$). \ Diện tích tăng thêm: $480-400=80$ (m$^2$).\par
$80:400=0{,}2=20\%$.
\dapso{$480$ m$^2$; diện tích tăng $20\%$.}}
\gap
\begin{probbox}{28}{\relax}
Sữa tươi chứa $4\%$ chất béo. Hỏi trong $3$ kg sữa tươi có bao nhiêu gam chất béo?
\end{probbox}
\probsol{
$3$ kg $=3\,000$ g. \ Lượng chất béo: $3\,000\times4:100=120$ (g).
\dapso{$120$ g.}}
\gap
\begin{probbox}{29}{\relax}
Một xe tải phải đi quãng đường $240$ km trong ba ngày. Ngày thứ nhất xe đi được $35\%$ quãng đường, ngày thứ hai đi được $40\%$ quãng đường. Hỏi ngày thứ ba xe phải đi bao nhiêu ki-lô-mét?
\end{probbox}
\probsol{
Ngày thứ ba đi: $100\%-35\%-40\%=25\%$ (quãng đường).\par
$240\times25:100=60$ (km).
\dapso{$60$ km.}}
\gap
\begin{probbox}{30}{\relax}
Quỹ lớp có $1\,500\,000$ đồng. Lớp dùng $40\%$ số tiền để mua sách, sau đó dùng $30\%$ số tiền còn lại để mua bút. Hỏi quỹ lớp còn lại bao nhiêu tiền?
\end{probbox}
\probsol{
Tiền mua sách: $1\,500\,000\times40:100=600\,000$ (đồng).\par
Còn lại: $1\,500\,000-600\,000=900\,000$ (đồng).\par
Tiền mua bút: $900\,000\times30:100=270\,000$ (đồng).\par
Quỹ lớp còn lại: $900\,000-270\,000=630\,000$ (đồng).
\dapso{$630\,000$ đồng.}}
\dividerbar{Thử thách -- Nâng cao}
\begin{probbox}{31}{\stars{1}}
Biết $25\%$ của một số là $18$. Tìm số đó.
\end{probbox}
\probsol{
$1\%$ của số đó là: $18:25=0{,}72$.\par
Số đó ($100\%$) là: $0{,}72\times100=72$.
\note{Hoặc: $25\%=\dfrac14$, nên số đó là $18\times4=72$.}
\dapso{$72$.}}
\gap
\begin{probbox}{32}{\stars{1}}
Một nhà máy sản xuất được $1\,380$ sản phẩm và đã vượt mức kế hoạch $15\%$. Hỏi theo kế hoạch, nhà máy phải sản xuất bao nhiêu sản phẩm?
\end{probbox}
\probsol{
Nhà máy đã thực hiện được: $100\%+15\%=115\%$ (kế hoạch).\par
$1\%$ kế hoạch là: $1\,380:115=12$ (sản phẩm).\par
Số sản phẩm theo kế hoạch: $12\times100=1\,200$ (sản phẩm).
\dapso{$1\,200$ sản phẩm.}}
\gap
\begin{probbox}{33}{\stars{1}}
Một cửa hàng bán một chiếc quạt với giá $240\,000$ đồng và được lãi $20\%$ giá vốn. Hỏi giá vốn của chiếc quạt là bao nhiêu?
\end{probbox}
\probsol{
Giá bán bằng: $100\%+20\%=120\%$ (giá vốn).\par
Giá vốn: $240\,000:120\times100=200\,000$ (đồng).
\dapso{$200\,000$ đồng.}}
\gap
\begin{probbox}{34}{\stars{1}}
Sau khi được giảm giá $25\%$, mẹ mua một chiếc áo khoác hết $360\,000$ đồng. Hỏi giá chiếc áo khi chưa giảm là bao nhiêu?
\end{probbox}
\probsol{
Mẹ đã trả: $100\%-25\%=75\%$ (giá chưa giảm).\par
Giá chưa giảm: $360\,000:75\times100=480\,000$ (đồng).
\dapso{$480\,000$ đồng.}}
\gap
\begin{probbox}{35}{\stars{1}}
Khi bán một món đồ chơi, cửa hàng được lãi $20\%$ giá \emph{bán}. Hỏi tiền lãi bằng bao nhiêu phần trăm giá \emph{vốn}?
\end{probbox}
\probsol{
Coi giá bán là $100\%$ thì tiền lãi là $20\%$, giá vốn là $100\%-20\%=80\%$ (giá bán).\par
Tỉ số phần trăm của tiền lãi so với giá vốn: $20:80=0{,}25=25\%$.
\dapso{$25\%$.}}
\gap
\begin{probbox}{36}{\stars{1}}
Một cửa hàng mua hàng với giá bằng $80\%$ giá niêm yết và bán ra với giá giảm $5\%$ so với giá niêm yết. Hỏi tiền lãi bằng bao nhiêu phần trăm giá vốn?
\end{probbox}
\probsol{
Coi giá niêm yết là $100\%$.\par
Giá vốn: $80\%$; \ giá bán: $100\%-5\%=95\%$; \ tiền lãi: $95\%-80\%=15\%$ (giá niêm yết).\par
Tỉ số phần trăm của tiền lãi so với giá vốn: $15:80=0{,}1875=18{,}75\%$.
\dapso{$18{,}75\%$.}}
\gap
\begin{probbox}{37}{\stars{1}}
Bác Dũng gửi tiết kiệm $10\,000\,000$ đồng với lãi suất $0{,}5\%$ một tháng. Tiền lãi tháng thứ nhất được cộng vào tiền gửi và tiếp tục tính lãi trong tháng thứ hai. Hỏi sau hai tháng, bác Dũng có tất cả bao nhiêu tiền?
\end{probbox}
\probsol{
Sau tháng thứ nhất: tiền lãi $10\,000\,000\times0{,}5:100=50\,000$ (đồng); \ có tất cả $10\,050\,000$ đồng.\par
Sau tháng thứ hai: tiền lãi $10\,050\,000\times0{,}5:100=50\,250$ (đồng); \ có tất cả $10\,050\,000+50\,250=10\,100\,250$ (đồng).
\dapso{$10\,100\,250$ đồng.}}
\gap
\begin{probbox}{38}{\stars{1}}
Một lọ chứa $300$ g nước muối có $6\%$ là muối. Cần đổ thêm bao nhiêu gam nước lọc để được nước muối có $4\%$ là muối?
\end{probbox}
\probsol{
Lượng muối: $300\times6:100=18$ (g); đổ thêm nước thì lượng muối không đổi.\par
Lượng nước muối mới: $18$ g bằng $4\%$ của nó, nên nặng $18:4\times100=450$ (g).\par
Cần đổ thêm: $450-300=150$ (g nước).
\dapso{$150$ g.}}
\gap
\begin{probbox}{39}{\stars{1}}
Trộn $200$ g nước muối có $10\%$ muối với $300$ g nước muối có $5\%$ muối. Hỏi muối chiếm bao nhiêu phần trăm hỗn hợp?
\end{probbox}
\probsol{
Muối trong lọ thứ nhất: $200\times10:100=20$ (g); \ trong lọ thứ hai: $300\times5:100=15$ (g).\par
Hỗn hợp nặng: $200+300=500$ (g), có $20+15=35$ (g) muối.\par
$35:500=0{,}07=7\%$.
\dapso{$7\%$.}}
\gap
\begin{probbox}{40}{\stars{1}}
Một lớp học có số học sinh nam chiếm $40\%$ số học sinh cả lớp. Sau khi $6$ học sinh nữ chuyển sang lớp khác thì số học sinh nam chiếm $50\%$ số học sinh của lớp. Hỏi lúc đầu lớp học có bao nhiêu học sinh?
\end{probbox}
\probsol{
Số học sinh nam không thay đổi.\par
Lúc đầu: nữ $60\%$, nam $40\%$, nên số nữ bằng $\dfrac{60}{40}=\dfrac32$ số nam.\par
Lúc sau: số nữ bằng số nam (mỗi bên $50\%$).\par
$6$ học sinh nữ chuyển đi ứng với $\dfrac32-1=\dfrac12$ số nam, nên số nam là $6\times2=12$ (học sinh).\par
Lúc đầu lớp có: $12:40\times100=30$ (học sinh).
\dapso{$30$ học sinh.}}
\gap
\begin{probbox}{41}{\stars{1}}
Giá một chiếc xe đạp được tăng thêm $25\%$. Hỏi phải giảm giá mới đi bao nhiêu phần trăm để trở lại giá cũ?
\end{probbox}
\probsol{
Coi giá cũ là $4$ phần; $25\%=\dfrac14$ nên giá mới là $5$ phần.\par
Muốn trở lại giá cũ phải bớt $1$ phần của giá mới: $1:5=0{,}2=20\%$.
\dapso{$20\%$.}}
\gap
\begin{probbox}{42}{\stars{1}}
Cửa hàng A bán một chiếc nồi cơm điện giá $200\,000$ đồng và giảm giá $15\%$. Cửa hàng B bán chiếc nồi cơm điện như thế giá $180\,000$ đồng và giảm giá $5\%$. Mua ở cửa hàng nào rẻ hơn và rẻ hơn bao nhiêu tiền?
\end{probbox}
\probsol{
Cửa hàng A: $200\,000\times85:100=170\,000$ (đồng).\par
Cửa hàng B: $180\,000\times95:100=171\,000$ (đồng).\par
Mua ở cửa hàng A rẻ hơn: $171\,000-170\,000=1\,000$ (đồng).
\dapso{cửa hàng A, rẻ hơn $1\,000$ đồng.}}
\gap
\begin{probbox}{43}{\stars{2}}
Một người bán hai chiếc bình, mỗi chiếc giá $240\,000$ đồng. Chiếc thứ nhất được lãi $20\%$ giá vốn, chiếc thứ hai bị lỗ $20\%$ giá vốn. Hỏi tính chung, người đó lãi hay lỗ, và lãi (lỗ) bao nhiêu tiền?
\end{probbox}
\probsol{
Chiếc thứ nhất: $240\,000$ đồng bằng $120\%$ giá vốn, nên giá vốn là $240\,000:120\times100=200\,000$ (đồng).\par
Chiếc thứ hai: $240\,000$ đồng bằng $100\%-20\%=80\%$ giá vốn, nên giá vốn là $240\,000:80\times100=300\,000$ (đồng).\par
Tổng giá vốn: $500\,000$ đồng; \ tổng tiền bán: $480\,000$ đồng.\par
Người đó lỗ: $500\,000-480\,000=20\,000$ (đồng).
\dapso{lỗ $20\,000$ đồng.}}
\gap
\begin{probbox}{44}{\stars{1}}
Một rạp hát giảm giá vé $20\%$ nên số người xem tăng thêm $30\%$. Hỏi số tiền bán vé của rạp thay đổi bao nhiêu phần trăm?
\end{probbox}
\probsol{
Giá vé mới bằng $80\%$ giá cũ; số người xem mới bằng $130\%$ số người cũ.\par
Tiền bán vé mới bằng: $80\%\times130\%=0{,}8\times1{,}3=1{,}04=104\%$ (tiền bán vé cũ).\par
Tiền bán vé tăng thêm: $104\%-100\%=4\%$.
\dapso{tăng $4\%$.}}
\gap
\begin{probbox}{45}{\stars{1}}
Cạnh của một hình vuông được tăng thêm $10\%$. Hỏi diện tích hình vuông tăng thêm bao nhiêu phần trăm?
\end{probbox}
\probsol{
Cạnh mới bằng $110\%$ cạnh cũ.\par
Diện tích mới bằng: $110\%\times110\%=1{,}1\times1{,}1=1{,}21=121\%$ (diện tích cũ).\par
Diện tích tăng thêm $21\%$.
\dapso{$21\%$.}}
\gap
\begin{probbox}{46}{\stars{1}}
Một xã có $5\,000$ người. Mỗi năm số dân của xã tăng thêm $2\%$ so với năm trước. Hỏi sau $2$ năm, xã đó có bao nhiêu người?
\end{probbox}
\probsol{
Sau $1$ năm: $5\,000+5\,000\times2:100=5\,000+100=5\,100$ (người).\par
Sau $2$ năm: $5\,100+5\,100\times2:100=5\,100+102=5\,202$ (người).
\dapso{$5\,202$ người.}}
\gap
\begin{probbox}{47}{\stars{2}}
Một khối đồng thau nặng $500$ g chứa $80\%$ là đồng, còn lại là thiếc. Cần thêm bao nhiêu gam đồng để khối mới chứa $90\%$ là đồng?
\end{probbox}
\probsol{
Thiếc chiếm $100\%-80\%=20\%$, nặng $500\times20:100=100$ (g). Thêm đồng thì lượng thiếc không đổi.\par
Trong khối mới, thiếc chiếm $100\%-90\%=10\%$, nên khối mới nặng $100:10\times100=1\,000$ (g).\par
Cần thêm: $1\,000-500=500$ (g đồng).
\dapso{$500$ g.}}
\gap
\begin{probbox}{48}{\stars{1}}
Một cửa hàng văn phòng phẩm khuyến mãi: ``Mua $4$ quyển vở, tặng $1$ quyển.'' Mua theo cách này thì giá mỗi quyển vở được giảm bao nhiêu phần trăm?
\end{probbox}
\probsol{
Trả tiền $4$ quyển mà được $5$ quyển, tức là tiết kiệm được tiền $1$ quyển trên giá của $5$ quyển.\par
$1:5=0{,}2=20\%$.
\dapso{$20\%$.}}
\gap
\begin{probbox}{49}{\stars{1}}
Một chị bán hàng mỗi tháng được nhận lương cơ bản $6\,000\,000$ đồng, cộng thêm $5\%$ giá trị số hàng bán được. Tháng trước chị nhận được $8\,500\,000$ đồng. Hỏi tháng trước chị đã bán được số hàng trị giá bao nhiêu tiền?
\end{probbox}
\probsol{
Tiền thưởng theo hàng bán: $8\,500\,000-6\,000\,000=2\,500\,000$ (đồng), bằng $5\%$ giá trị hàng bán được.\par
Giá trị số hàng: $2\,500\,000:5\times100=50\,000\,000$ (đồng).
\dapso{$50\,000\,000$ đồng.}}
\gap
\begin{probbox}{50}{\stars{2}}
Trong một cuộc thi Toán, $60\%$ số học sinh dự thi đến từ thành phố, còn lại đến từ nông thôn. Có $25\%$ số học sinh thành phố và $40\%$ số học sinh nông thôn đạt giải. Hỏi số học sinh đạt giải chiếm bao nhiêu phần trăm số học sinh dự thi?
\end{probbox}
\probsol{
Học sinh nông thôn chiếm: $100\%-60\%=40\%$ (số học sinh dự thi).\par
Học sinh thành phố đạt giải: $60\%\times25\%=0{,}6\times0{,}25=0{,}15=15\%$ (số học sinh dự thi).\par
Học sinh nông thôn đạt giải: $40\%\times40\%=0{,}4\times0{,}4=0{,}16=16\%$ (số học sinh dự thi).\par
Tổng số học sinh đạt giải: $15\%+16\%=31\%$.
\dapso{$31\%$.}}
